\section{Finite-horizon decision equivalence and a sharp causal state}\label{sec:decision-horizon}
We now connect the attainable class to the information needed for a
specified future experiment. The parameter interval is affinely identified
with $[0,1]$. Fix a full-support probability $\pi_0$ there. The apparatus
has a finite raw alphabet, positive polynomial probabilities of degree at
most $q\ge1$, and a fixed finite lookup gate class. Assume its attainable
failure polynomials contain a nonempty open subset of $\mathcal P_q$.
The two apparatuses of Theorem~\ref{thm:calibrated-rank} satisfy these
hypotheses with $q=3$ and $q=2$. Throughout this section $0<\eta<1/2$ is fixed, and gate values range over
the full finite-dimensional lookup cube $[\eta,1-\eta]^{|\mathcal X|}$.
All known commands belong to the history;
they are not an uncharged external memory available after compression.

After $n$ reports let the normalized posterior density relative to $\pi_0$
be $P$, so $P\in\mathcal P_{qn}$, $P>0$ and $\pi_0P=1$. This notation
normalizes the likelihood by its prior integral, not by a coefficient sum.
For $m$ remaining cartridges, two such densities are called
\emph{operationally equivalent} if every common admissible continuation
policy gives the same future report law. Independent policy randomization
is allowed. Equality therefore gives the same expected value for every
bounded measurable payoff of future reports, commands and a terminal
prediction or decision. A payoff may not in this definition inspect the
hidden parameter directly. The future task specification is the same for
the two histories. An externally prescribed payoff depending on a past
observation requires its own task state and is not erased by this definition. The unrestricted terminal payoffs $G(R,d)$ of
Section~\ref{sec:control} retain the full-posterior state; they are not
silently included in this smaller quotient. All state minimality assertions
refer to one common representation on the admitted known-command histories.
They are not a lower bound for a code customized to one fixed finite list
of possible prefixes.

\begin{lemma}[Polynomial span of physically executed probes]\label{lem:probe-span}
There is a finite list of $m$-cartridge nonadaptive gate sequences whose
all-failure probabilities $H_1,\ldots,H_s$ span $\mathcal P_{qm}$. Every
$H_j$ is a product of $m$ attainable positive failure factors and is the
probability of an event in an actually executed experiment. For the shared
menu in Section~\ref{sec:calibrated}, all $4^m$ words are such a list for
both apparatuses, with their respective values of $q$.
\end{lemma}
\begin{proof}
An open subset of $\mathcal P_q$ spans $\mathcal P_q$, so select finitely
many attainable factors spanning it. Their $m$-fold products span every
product of $m$ members of $\mathcal P_q$. Each monomial $t^k$, $0\le k\le qm$,
is such a product: distribute $k$ among $m$ integers between zero and $q$.
Thus the span is exactly $\mathcal P_{qm}$. The probability formula follows
by independence of fresh cartridges conditional on the same parameter;
all $m$ cartridges are consumed even if an earlier report is not a failure.
The shared four-factor menu has the stated one-step spans, and the same
argument applies. One may select $qm+1$ independent probes for a single
apparatus, but the shared ablation retains the same larger index set.
\end{proof}

\begin{theorem}[Finite-horizon decision quotient]\label{thm:decision-quotient}
For histories after $n$ cartridges, the following are equivalent:
\begin{enumerate}
\item their posteriors are operationally equivalent for $m$ future cartridges;
\item $\int t^jP(t)\pi_0(dt)$ agrees for $j=1,\ldots,qm$;
\item all the probe predictions $p_j(P)=\int H_jP\,d\pi_0$ agree.
\end{enumerate}
The attainable prediction set lies in an affine space of dimension
\begin{equation}\label{eq:decision-dimension}
 d_{n,m}=\min\{qn,qm\}=q\min\{n,m\},
\end{equation}
and contains a relatively open set in that space, with a continuous local
section to histories consisting solely of failures. Consequently its
smallest exact continuous Euclidean state dimension is $d_{n,m}$.
A state of this dimension represents all attainable histories, not merely
the local family used to prove sharpness.
\end{theorem}
\begin{proof}
For a fixed realized future word a common feedback policy chooses known
commands independent of the unknown parameter. Its conditional likelihood
is a product of at most $m$ polynomials of degree at most $q$, hence belongs
to $\mathcal P_{qm}$. Equality of the listed moments gives equality of every
word probability and thus the full future law. Integrating a common
independent random seed gives the same statement for randomized policies.
Conversely operational equivalence includes each nonadaptive probe and its
all-failure event; Lemma~\ref{lem:probe-span} then gives equality of all
moments. This proves the equivalences, including arbitrary feedback in the
finite lookup class, rather than only the finite probe menu.

Set $a=qn$, $b=qm$, and let
$E_a=\{Q\in\mathcal P_a:\pi_0Q=0\}$. The posterior family lies in
$1+E_a$. On this affine space the prediction map is linear up to a fixed
translation. Its tangent map is
\[
 L:E_a\longrightarrow\mathbb R^s,\qquad
 LQ=(\pi_0(H_jQ))_{j=1}^s.
\]
It has rank $\min\{a,b\}$. If $b\ge a$, a kernel vector is orthogonal to
all of $\mathcal P_b$, hence to itself; full support implies $Q=0$.
If $b<a$, restrict the domain to $E_b\subset E_a$. The same argument
shows injectivity there, so the rank is at least $b$. It is at most $b$
since the tests span $\mathcal P_b$ and the constant test annihilates
$E_a$. Full support is used through positive definiteness of the polynomial
Gram form $\pi_0(Q^2)$, not through an assumed density of the prior.

The open attainable factor family contains $n$ exact-degree, pairwise
coprime factors. The product differential in
Theorem~\ref{thm:rank-deficient}, now with $W=\mathcal P_q$, has rank
$qn+1$ before normalization and $qn$ after normalization by $\pi_0$.
This gives an open subset of $1+E_a$ with a smooth local section to gate
histories. Its image under the rank-$d_{n,m}$ linear map $L$ is relatively
open; restrict a linear right inverse to a small ball and compose the two
sections. All those histories have probability at least $\eta^n$ in the
bounded gate class, at every parameter. They are retained, not discarded
by conditioning away the physical preparation.

Any exact continuous state composed with this section is a continuous
injection of an open $d_{n,m}$-ball into its Euclidean state space, and
invariance of domain proves the lower bound. For the upper bound use
normalized polynomial coordinates when $a\le b$, and the first $b$ moments
when $b<a$. Each is continuous on all positive-likelihood histories and
has exactly $d_{n,m}$ free coordinates. Equivalently use coordinates in
the affine hull of the prediction vector. This proves the global upper
bound as well as local sharpness. The cases $n=0$ or $m=0$ have a single
operational class and dimension zero.
\end{proof}

\begin{theorem}[An exact causal state with a tent-shaped memory profile]\label{thm:causal-horizon}
For a prescribed total budget $N$ and observable future-payoff problems,
there is a continuous exact causal state at time $n$ with
$q\min\{n,N-n\}$ free real coordinates. Its maximum is
$q\lfloor N/2\rfloor$. With the prior moments through degree $qN$
precomputed as read-only apparatus data, all updates use finite polynomial
arithmetic. This is a statement about the posterior-dependent state;
it does not charge immutable calibration tables or count global policy
optimization as a filter update.
\end{theorem}
\begin{proof}
While $n\le N-n$, store the normalized degree-$qn$ polynomial. A reported
likelihood factor $f(t)=\sum_{i=0}^q f_it^i$ is multiplied into it and
normalized by its prior integral. Once $n>N-n$, store the moments
$M_j=\pi(t^j)$, $1\le j\le q(N-n)$, with $M_0=1$ known. At the changeover
these moments are computed exactly from the polynomial and prior moments.
If $m=N-n$ and another report is observed, the required smaller list is
\begin{equation}\label{eq:shrinking-update}
 M'_j=\frac{\sum_{i=0}^q f_iM_{i+j}}
             {\sum_{i=0}^q f_iM_i},\qquad 0\le j\le q(m-1).
\end{equation}
Only old moments of order at most $qm$ appear. The denominator is positive
on the actual record; assign a fixed value on a common null branch if
endpoint gates are admitted. Theorem~\ref{thm:decision-quotient} proves
sufficiency after every update. The changeover uses prior moments at most
$qn+q(N-n)=qN$. If desired, fix a normalization functional nonzero on every
positive factor and eliminate one polynomial coefficient explicitly; no
extra free coordinate is hidden in the coefficient array. The dimension
lower bound already holds for each checkpoint's all-failure family.
\end{proof}

\paragraph{What changes when a terminal parameter loss is prescribed.}
The operational quotient above is not a claim that a few moments determine
all losses of the hidden radius. A useful precise extension is available:
if $G(t,d)$ is polynomial of degree at most $s_0$, with bounded measurable
coefficients in $d$, then all future expected payoffs obtained by multiplying
$G$ by bounded observable path rewards depend only on moments through
$qm+s_0$. This follows by multiplying the future word polynomial by $G$
and integrating. A matching lower bound requires that the specified
terminal-decision class actually spans the additional tests; it is not
inferred from its nominal degree. For an arbitrary continuous $G$, the
full likelihood state and the original Bellman theorem remain in force.

\paragraph{Relation to predictive representations.}
Prediction-based states and action-conditional tests are established ideas
\cite{LittmanSutton,SinghJamesRudary}. The argument here does not rename that
principle. It identifies the attainable posterior manifold, computes its
finite-horizon observable quotient exactly, and supplies a causal realization
and a costed finite-memory decision experiment. The dimension in
\eqref{eq:decision-dimension} is an intrinsic continuous dimension at a
fixed checkpoint, not a horizon-independent linear rank of an infinite
system-dynamics matrix.
